\begin{lemma}[An affine envelope above the forced midpoint]\label{fprof:affine}
Suppose \(g\) is 1-Lipschitz on \([0,N]\), \(g(0)=0\), and
$$
 ax\le g(x)\le bx,\qquad 0<a\le b\le\tfrac12.
$$
For \(N/2\le n_0<N\), there is an admissible decreasing chain from \(n_0\) through
\(N/2\) to zero, with \(O(1+\log(N/(N-n_0)))\) edges and cost at most
\begin{equation}\label{fprof:affinecost}
 \left(\frac{1+2b-4a}{8}+\frac{b-a}{2(1+a)}\right)N.
\end{equation}
On a grid of mesh \(T\) containing \(N/2,3N/4,n_0\), add \(O(T)\).
If the barriers hold within an additive error \(eN\), add \(O(KeN+KT)\),
where \(K\) bounds the number of edges.
\end{lemma}
\begin{proof}
First assume a grid profile. Put \(l=N/2\), \(r=3N/4\), and
$$
 \alpha=\frac{b-1/2}{2}N,\qquad
 L(x)=x-\frac N2+\alpha-g(x),\qquad V(x)=\frac{L(x)_+}{2}.
$$
For \(x\ge l\), the upper barrier gives
$$
 g(x)\le bx\le \tfrac12x+\alpha.
$$
The function \(L\) is nondecreasing. Almost everywhere on \(\{L>0\}\),
$$
 (-g')_+\le\frac{1-g'}2=V'.
$$
For \(g'\ge0\) the right side is nonnegative; for \(g'<0\) the inequality
is equivalent to \(g'\ge-1\). Also \(c_g(m,n)\le\int_m^n(-g')_+\).

If \(n_0\le r\), proceed directly to the last paragraph of the chain construction.
Otherwise let \(q\) be the largest grid point in \([l,N]\) with \(L(q)\le0\).
It exists since \(L(l)=\alpha-g(l)<0\). If \(n_0>\max(r,q+T)\), repeatedly
replace the current \(n\) by \(\max(r,q+T,2n-N)\) until this target is reached.
All edges are admissible and lie where \(L>0\), so their total cost is at most
\(V(n_0)\). Apart from the last edge, \(N-n\) doubles.
If necessary, pass from \(q+T\) to \(q\), at cost at most \(T\).
This step is admissible because it is used only with \(q+T\le n_0<N\), hence
\(q+2T\le N\). We have either reached \(r\), or a point where \(L\le0\).

If the current point \(n>r\) satisfies \(L(n)\le0\), set \(m_0=2n-N>l\).
The affine envelope gives
$$
 g(m_0)\le\tfrac12m_0+\alpha
 =n-\tfrac N2+\alpha\le g(n).
$$
Choose the leftmost grid minimizer \(m\) on \([m_0,n]\).
It is strictly smaller than \(n\), and the resulting admissible edge has zero cost.
Continue until the current point belongs to \([l,r]\).
Monotonicity of \(L\) keeps this continuation in \(\{L\le0\}\).
Merge consecutive zero-cost edges greedily whenever their union is admissible.
The union remains zero-cost. If two consecutive resulting edges cannot be merged,
their combined complementary depths increase by a factor greater than two.
Thus this part has \(O(1+\log(N/(N-n_0)))\) edges.

From a point \(n\in[l,r]\), jump to \(l\); this edge is admissible.
Writing \(v=g(l)\), for every \(x\ge l\) one has
$$
 g(x)\ge\max\{ax,v-(x-l)\}\ge\frac{a(v+l)}{1+a}.
$$
The final inequality follows by minimizing the maximum of these two affine functions.
Consequently the last edge costs at most
$$
 \frac{v-al}{1+a}\le\frac{(b-a)N}{2(1+a)}.
$$
The edge from \(l\) to zero has zero cost. Finally,
$$
 V(n_0)\le V(N)\le\frac{1+2b-4a}{8}N,
$$
where the last numerator is nonnegative since \(a\le b\le1/2\).
This proves the grid claim and the length bound.

For a continuous profile, interpolate on finer grids containing \(0,l,r,N\).
Move \(n_0\) down by at most one mesh length first if necessary; for sufficiently
fine meshes this extra edge is admissible and costs at most that length.
Interpolation preserves the Lipschitz constant and barriers. The length bound is uniform.
Pad the endpoint vectors by repetitions and extract a convergent subsequence.
Uniform convergence of profiles and convergence of endpoints imply convergence of interval
minima. Admissibility is closed; discard repeated endpoints. This proves the exact real bound.
For approximate barriers, clamp the profile between \(ax\) and \(bx\), then interpolate
grid values. The uniform change is \(O(eN+T)\); minimum and maximum preserve the
Lipschitz constant. Each edge cost changes by at most twice this uniform error.
\end{proof}

Relative to the preceding midpoint lemma with \(b=1/2\) and actual upper slope
\(b_0\le1/2\), the gain is \((1/2-b_0)N/4\).
It comes from using the upper bound only above the mandatory midpoint.
The bound is not asserted to be optimal even for this chain problem.
